\documentclass[institutional,screen,none]{../cls/ua-engineering-v6-beamer}
% !TeX encoding = UTF-8
\UASetUniversity{Universidad Austral}
\UASetFaculty{Facultad de Ingeniería}
\UASetDepartment{Departamento de Computación}
\UASetCampus{Pilar}
\UASetCareer{Ingeniería Informática}
\UASetCommission{Comisión A}
\UASetProfessor{Equipo docente}
\UASetSemester{Segundo cuatrimestre 2026}
\UASetCourse{Sistemas Distribuidos}
\UASetDate{2026}


\UASetClassNumber{14}
\UASetClassTitle{Incidentes y cultura de confiabilidad}
\title{\UACourse}
\subtitle{Clase 14 --- Incidentes, postmortems, aprendizaje organizacional y cultura de confiabilidad}

\begin{document}

{
\setbeamertemplate{footline}{}
\begin{frame}[plain]
\titlepage
\end{frame}
}

\UASectionSlide{El incidente como hecho normal}

\begin{frame}{La intuición equivocada}
Un incidente suele vivirse como:
\begin{itemize}
\item excepción vergonzosa
\item error puntual
\item fracaso individual
\end{itemize}

\pause

\begin{alertblock}{Idea correcta}
En sistemas complejos, los incidentes son eventos esperables y su valor depende de cuánto aprende la organización de ellos.
\end{alertblock}
\end{frame}

\begin{frame}{Qué es un incidente}
Un incidente es un evento que degrada o amenaza:
\begin{itemize}
\item disponibilidad
\item latencia
\item corrección
\item seguridad
\item cumplimiento de SLOs
\end{itemize}
\end{frame}

\UASectionSlide{Gestión de incidentes}

\begin{frame}{Fases}
\begin{itemize}
\item detección
\item clasificación
\item respuesta
\item mitigación
\item comunicación
\item recuperación
\item aprendizaje posterior
\end{itemize}
\end{frame}

\begin{frame}{Prioridad}
No toda falla tiene la misma urgencia.

\pause

Se prioriza según:
\begin{itemize}
\item impacto al usuario
\item alcance
\item duración
\item consumo de error budget
\item riesgo de escalamiento
\end{itemize}
\end{frame}

\begin{frame}{Incident Command}
Buenas prácticas:
\begin{itemize}
\item alguien coordina
\item alguien investiga
\item alguien comunica
\item alguien registra timeline
\end{itemize}

\pause

\begin{block}{Clave}
Durante el incidente, coordinación y claridad importan más que brillantez individual.
\end{block}
\end{frame}

\UASectionSlide{Postmortems}

\begin{frame}{Postmortem}
\begin{block}{}
Documento y proceso de análisis posterior al incidente para entender qué pasó, por qué pasó y cómo reducir repetición.
\end{block}
\end{frame}

\begin{frame}{Qué debe incluir}
\begin{itemize}
\item impacto
\item timeline
\item detección
\item causa raíz
\item factores contribuyentes
\item mitigación
\item acciones de mejora
\end{itemize}
\end{frame}

\begin{frame}{Blameless}
\begin{itemize}
\item no significa ausencia de responsabilidad
\item significa evitar castigo como sustituto del análisis
\end{itemize}

\pause

\begin{alertblock}{Clave}
El objetivo del postmortem no es encontrar culpables, sino aumentar la capacidad de la organización.
\end{alertblock}
\end{frame}

\UASectionSlide{Root Cause y factores contribuyentes}

\begin{frame}{Causa raíz}
Un incidente rara vez tiene una sola causa.

\pause

Suele haber:
\begin{itemize}
\item disparador inmediato
\item condiciones habilitantes
\item fallas de detección
\item fallas de mitigación
\end{itemize}
\end{frame}

\begin{frame}{Ejemplo}
\begin{itemize}
\item DB lenta
\item retries mal calibrados
\item breaker ausente
\item alertas tardías
\item comunicación confusa
\end{itemize}

\pause

\begin{block}{Lección}
El incidente real es una cadena, no un único bug.
\end{block}
\end{frame}

\UASectionSlide{Aprendizaje organizacional}

\begin{frame}{Aprender del incidente}
Un buen postmortem produce:
\begin{itemize}
\item fixes técnicos
\item mejoras de observabilidad
\item cambios en runbooks
\item ajustes en alertas
\item cambios en procesos
\end{itemize}
\end{frame}

\begin{frame}{Anti-patrones}
\begin{itemize}
\item cerrar incidente sin aprender
\item hacer postmortem solo para “cumplir”
\item listar acciones sin owner ni fecha
\item reducir todo a error humano
\end{itemize}
\end{frame}

\UASectionSlide{Cultura de confiabilidad}

\begin{frame}{Qué cultura queremos}
\begin{itemize}
\item incidentes visibles, no ocultos
\item aprendizaje explícito
\item responsabilidad sin cultura de miedo
\item mejora continua
\end{itemize}
\end{frame}

\begin{frame}{Idea central}
\begin{block}{}
La confiabilidad no es solo propiedad del sistema; es también propiedad de la organización que lo opera.
\end{block}
\end{frame}

\UASectionSlide{Cierre}

\begin{frame}{Idea final}
\begin{block}{}
Una organización madura no se define por no tener incidentes, sino por responder, aprender y mejorar de manera sistemática cuando los tiene.
\end{block}
\end{frame}

\end{document}
